\scanpage{54}
the objects of $\unsure{\bar C|\widetilde C}$ \deleted{are} sheaves (or only separated presheaves)

In fact, we get

\uline{Th. 3} (Giraud) $T\longrightarrow\widetilde C_T$ is a bijective correspondance between topologies on $C$ and \inserted{strictly full} subcategories \inserted{$E$} of $\widehat C$ s.th. the inclusion functor $i\colon E\longrightarrow\widehat C$ has a left adjoint which is left exact. It is order reversing (\deleted{\unsure{finer}} $T\leq T'\Longleftrightarrow\widetilde C_{T'}\subset\widetilde C_T$).

\uline{Th. 5} Let $E$ be a category. Then $E\simeq\widetilde C$ some (small) $C$ iff $E$ satisfies a) b) c) d).

\deleted{We can take for $C$ the class} \inserted{We may always choose} \deleted{of such a} $C$ with a topology \deleted{\unsure{finer}} \inserted{less fine} than the canonical one (ie. $C\subset\widetilde C$). The choice \deleted{of such a} \inserted{for \deleted{\unsure{a}} a given category} $C$ \inserted{of a top. $T$ on $C$} \deleted{and \unsure{equivce}} \inserted{finer than the can. one, and an \unsure{équiv.}} $\widetilde C_T\overset{\varphi}{\simeq}E$ is equivalent \deleted{(up to equivalence for $\varphi$, see \unsure{p.})} to a \deleted{full subcategory of $E$ which is generating, and the topology} \inserted{fully faithful} functor $C\longrightarrow E$ such that the image of $C$ is generating $E$; the topology \inserted{$T$} of $C$ is \uline{induced} by the canonical topology of $E$, via the given functor....

Continuous functor between sites, induced topology, and:

\inserted{(for instance by $C\xrightarrow{\varepsilon_C}\deleted{E}=\widetilde C$)}

% large return arrow points upward from Th. 4 to Th. 5
\uline{Th. 4} \uline{Comparison lemma} $C\underset{i}{\lhook\joinrel\longrightarrow}C'$ \inserted{$U$-cat.} inclusion de sites, $i'$ \unsure{plein\textsuperscript{mt}|pleinemt|pleinement} fidèle, top. de $C$ induite.

Si tout objet de $C'$ peut être couvert par des objets de $C$, alors $\widetilde C'\xrightarrow{\simeq}\widetilde C$.

(il y a réciproque si $C$ est un $U$-site dont la top est moins fine que la top. can.)

\margin{Analogue bien connu en top. \unsure{ordinaire|relativem\textsuperscript{t}} possibilité de définir des $\widetilde C$ équivalents de bien des façons différentes.}
% archive sheet number: 53
